% ============================================================================
% Fichier  : partie5/P05_C16_antiforensique.tex : FORN (LPI342)
% Sources : ISO/IEC 27042 ; NIST SP 800-86. Opération TSANAGA.
% ============================================================================

\chapter{L'Anti-Forensique et le Chiffrement}
\label{chap:anti}
\chaptermark{L'anti-forensique}

\epigraph{%
    « Effacer une trace en laisse souvent une autre~:\\
    celle de l'effacement. »%
}{Principe du cours}

\niveauChapitre{Expert}{Partie~IV.}

\begin{boxobjectifs}
\begin{itemize}
    \item Reconnaître les techniques d'anti-forensique.
    \item Comprendre l'effacement sécurisé et ses limites.
    \item Traiter le chiffrement dans une enquête.
    \item Repérer les traces d'une dissimulation.
    \item Rester honnête sur ce que l'on ne peut établir (NEK).
\end{itemize}
\end{boxobjectifs}

\bigskip

% ============================================================================
\section*{Question directrice : si M. NDONGO a effacé et chiffré ses traces, l'enquête est-elle bloquée ?}
% ============================================================================

Un suspect averti efface, chiffre, dissimule. Ces techniques compliquent
l'enquête mais laissent elles-mêmes des traces. Ce chapitre apprend à les
reconnaître et à dire honnêtement ce qui reste hors d'atteinte.

\medskip
\textbf{Plan du chapitre.}
\begin{enumerate}
    \item Les techniques d'anti-forensique (section~\ref{sec:c16-techniques}).
    \item L'effacement sécurisé (section~\ref{sec:c16-effacement}).
    \item Le chiffrement (section~\ref{sec:c16-chiffrement}).
    \item L'honnêteté du NEK (section~\ref{sec:c16-nek}).
\end{enumerate}

\separateur

% ============================================================================
\section{Les Techniques d'Anti-Forensique}
\label{sec:c16-techniques}
% ============================================================================

\begin{tableaucours}{L{3.6cm}L{8.6cm}}{%
    \tableauHeader{Technique} &
    \tableauHeader{Trace qu'elle laisse souvent}
}
Effacement sécurisé & Absence anormale, outil installé, date d'exécution \\
Chiffrement & Conteneur, extension, outil, clé en mémoire \\
Modification d'horodatage & Incohérences entre artefacts \\
Dissimulation (renommage) & Signature qui trahit le vrai type \\
Effacement des journaux & Événement de nettoyage, trou dans la suite \\
\end{tableaucours}
\vspace{-0.8em}
\begin{center}{\small\itshape Tableau~16.1. Techniques d'anti-forensique et leurs traces}\end{center}

% ============================================================================
\section{L'Effacement Sécurisé}
\label{sec:c16-effacement}
% ============================================================================

L'\termecle{effacement sécurisé} réécrit l'espace pour empêcher la
récupération. Mais il laisse des indices~: l'outil utilisé, sa date
d'exécution, l'absence anormale de données là où il devrait y en avoir. Un
effacement ciblé, juste avant une saisie, est en soi un élément.

% ============================================================================
\section{Le Chiffrement}
\label{sec:c16-chiffrement}
% ============================================================================

\begin{boxalert}
\textbf{Violation typique.} Tenter de « casser » un chiffrement fort par force
brute, sans espoir réaliste.

\smallskip
\textbf{Règle~:} contre un chiffrement fort, on cherche la clé ailleurs
(mémoire capturée à chaud, mot de passe noté, phrase de passe faible), pas la
force brute. Dans TSANAGA, la clé était en mémoire
(Chapitre~\ref{chap:memoire}).
\end{boxalert}

% ============================================================================
\section{L'Honnêteté du NEK}
\label{sec:c16-nek}
% ============================================================================

\begin{boxscenario}{Opération TSANAGA. Ce qui reste hors d'atteinte}
\textbf{Constatation.} Un second conteneur chiffré, non monté à la saisie,
reste inaccessible~: aucune clé trouvée, chiffrement fort.

\medskip
\textbf{Analyse} (confiance élevée sur l'impossibilité). On ne peut établir le
contenu de ce conteneur~: c'est un \sigle{NEK}, à déclarer tel quel. On ne
suppose pas son contenu~; on rapporte qu'il existe et reste chiffré. L'honnêteté
sur les limites renforce la crédibilité du reste.
\end{boxscenario}

\begin{boxCRO}
\textbf{Déclarer l'inaccessible.} Prétendre connaître le contenu d'un
conteneur non ouvert dégraderait la \cro{0}{-}{-}~: affirmation non fiable,
non opposable. Le NEK préserve la valeur de l'ensemble du rapport.
\end{boxCRO}

\begin{boxmetier}{Expert judiciaire}
\textbf{Dire ce qu'on ne sait pas.} L'expert qui déclare honnêtement un NEK
(« ce conteneur reste chiffré, son contenu est indéterminé ») est plus crédible
que celui qui comble les vides par des suppositions. La limite assumée est une
force.
\end{boxmetier}

\separateur

\begin{boxresume}
\textbf{Ce qu'il faut retenir de ce chapitre~:}
\begin{itemize}
    \item \textbf{L'anti-forensique} laisse souvent la trace de
          l'effacement.
    \item \textbf{L'effacement sécurisé} réécrit, mais laisse des indices.
    \item \textbf{Contre le chiffrement fort}, on cherche la clé ailleurs.
    \item \textbf{Ce qui reste inaccessible} est un NEK, à déclarer.
    \item \textbf{L'honnêteté sur les limites} renforce la crédibilité.
\end{itemize}
\end{boxresume}

\begin{boxinfo}
\textbf{Points de contrôle de l'élève.} Je sais~: reconnaître les techniques
d'anti-forensique~; expliquer l'effacement sécurisé~; traiter le
chiffrement~; déclarer un NEK~; défendre l'honnêteté des limites.
\end{boxinfo}

\bigskip

\begin{boxexo}[title={\ding{45}\quad Exercice~16.1. Techniques \niveaubase}]
Citez trois techniques d'anti-forensique.
\end{boxexo}

\begin{boxexo}[title={\ding{45}\quad Exercice~16.2. Effacement \niveaubase}]
Quelles traces laisse un effacement sécurisé~?
\end{boxexo}

\begin{boxexo}[title={\ding{45}\quad Exercice~16.3. Chiffrement \niveaubase}]
Contre un chiffrement fort, où cherche-t-on la clé~?
\end{boxexo}

\begin{boxexo}[title={\ding{45}\quad Exercice~16.4. L'effacement comme indice \niveaumoyen}]
Pourquoi un effacement juste avant une saisie est-il un élément~?
\end{boxexo}

\begin{boxexo}[title={\ding{45}\quad Exercice~16.5. NEK \niveaumoyen}]
Que faire du second conteneur chiffré inaccessible~?
\end{boxexo}

\begin{boxexo}[title={\ding{45}\quad Exercice~16.6. Force brute \niveaumoyen}]
Pourquoi la force brute contre un chiffrement fort est-elle vaine~?
\end{boxexo}

\begin{boxexo}[title={\ding{45}\quad Exercice~16.7. CRO \niveaumoyen}]
Pourquoi supposer le contenu d'un conteneur dégrade-t-il la preuve~?
\end{boxexo}

\begin{boxexo}[title={\ding{45}\quad Exercice~16.8. Honnêteté \niveaumoyen}]
En quoi déclarer un NEK renforce-t-il la crédibilité~?
\end{boxexo}

\begin{boxexo}[title={\ding{45}\quad Exercice~16.9. Détecter \niveauexpert}]
Comment détecter une modification d'horodatage dans TSANAGA~?
\end{boxexo}

\begin{boxexo}[title={\ding{45}\quad Exercice~16.10. Note de limites \niveauexpert}]
Rédigez la section « limites de l'expertise » du rapport TSANAGA.
\end{boxexo}

\bigskip

\begin{boxinfo}
\textbf{Pour aller plus loin.}
\begin{itemize}
    \item ISO/IEC 27042~\cite{iso27042}~; NIST SP~800-86~\cite{nist80086}.
\end{itemize}
\end{boxinfo}

\begin{boxinfo}
\textbf{Transition.} Il faut maintenant assembler les domaines en une
histoire~: le Chapitre~\ref{chap:reconstitution} traite de la corrélation et
de la reconstitution chronologique.
\end{boxinfo}
